%==============================================================================
\chapter{Contraintes et exigences techniques}
%==============================================================================

\section{Architecture}
Architecture \textbf{scalable, évolutive, flexible, multi-niveaux et faiblement
couplée}, s'appuyant sur \textbf{Spring Boot 3} (Java 17 LTS) comme
infrastructure principale, avec une séparation stricte entre l'API de service et
le client de présentation.

\section{Exigences techniques retenues}
\begin{xltabular}{\textwidth}{>{\bfseries}p{4cm} X}
\toprule
\textbf{Exigence} & \textbf{Description} \\
\midrule
\endhead
Architecture en couches & Séparation contrôleur / service / persistance, inversion de dépendance assurée par le conteneur Spring. \\
API REST & Contrôleurs Spring MVC exposant une API versionnée, décrite par un contrat \textbf{OpenAPI 3}. \\
Présentation & Client \textbf{React} en \textbf{TypeScript}, découplé du serveur et consommant l'API REST. \\
Services métier & Composants Spring (\texttt{@Service}) encapsulant les règles de gestion. \\
Persistance & \textbf{Spring Data JPA} pour le CRUD, \textbf{jOOQ} pour les requêtes analytiques et géospatiales. \\
Internationalisation & Textes de l'IHM externalisés en fichiers de ressources par locale, ouverture à la traduction hors-ligne sans limite de langues, FR, EN et AR au minimum en démonstration. \\
Configuration technique & Paramétrage par profils Spring (\texttt{application.yml}) et variables d'environnement : sources de données, points d'accès, secrets. \\
Configuration métier & Seuils, unités et modules actifs dans le fichier texte \texttt{ConfigZumm.ini}, ajustable par l'exploitant sans recompilation ni redéploiement. \\
Sécurité inter-composants & Certificats \textbf{X.509} et protocole \textbf{TLS 1.3}. \\
Authentification utilisateur & \textbf{OpenID Connect} via \textbf{Keycloak}, avec fédération d'identité Google et jetons \textbf{JWT}. \\
API tierce       & Point d'entrée REST exposant l'opération \texttt{getZummHoneyActualQuantity(int zummID)} pour récupérer la quantité de miel d'une ruche, publiée au contrat OpenAPI. \\
Front-end        & JavaScript / TypeScript et bibliothèques libres, dans le respect de l'ergonomie imposée. \\
\bottomrule
\end{xltabular}

\begin{encadre}
\textbf{Stabilité des identifiants :} le changement de socle technique ne
modifie ni le domaine métier ni les identifiants publics. L'opération
\texttt{getZummHoneyActualQuantity} conserve son nom et sa signature, quelle que
soit la technologie de transport retenue pour l'exposer.
\end{encadre}

\section{Exigences non fonctionnelles}
\begin{itemize}
  \item \textbf{Ergonomie, convivialité et simplicité} d'utilisation, consistance de l'interface et qualité du design infographique (importance capitale).
  \item \textbf{Évolutivité} : ouverture à l'ajout de langues, d'unités et d'indicateurs.
  \item \textbf{Interopérabilité} : API interrogeable par des applications tierces à partir d'un contrat publié.
  \item \textbf{Sécurité} : chiffrement des échanges et authentification forte déléguée à un fournisseur d'identité.
  \item \textbf{Qualité de conception} : respect des principes \textbf{SOLID} et recours aux \textbf{design patterns} (patrons de conception) pour garantir un couplage faible et l'évolutivité, comme détaillé en annexe~\ref{chap:solid}.
\end{itemize}

\begin{encadre}
\textbf{Principes de conception retenus :} l'architecture applique les cinq
principes \textbf{SOLID} (responsabilité unique, ouvert/fermé, substitution de
Liskov, ségrégation des interfaces, inversion de dépendance) et un ensemble de
\textbf{design patterns} (Composite, State, Strategy, Observer, Command, Factory
Method, Builder, Adapter, Facade, Singleton, Proxy). Leur application au
système, justifiée « avant / après », est détaillée dans
l'annexe~\ref{chap:solid}.
\end{encadre}

\section{Questions de conception à couvrir}
La conception devra répondre explicitement aux questions suivantes :
\begin{enumerate}
  \item Quels sont les \textbf{inputs} (liste, dictionnaire de données et types) et \textbf{outputs} (liste, dictionnaire, types, formules et requêtes de calcul) du système ?
  \item De quoi se compose la solution ?
  \item Quels sont les utilisateurs types et leurs droits d'accès aux fonctionnalités ?
  \item Comment le système réalise-t-il les opérations utilisateur ?
  \item Comment le système est-il packagé et déployé ?
  \item Comment les éléments interagissent-ils et dans quel ordre / chronologie ?
  \item Pourquoi, comment et où telle technologie a-t-elle été utilisée ?
\end{enumerate}

\section{Justification des choix technologiques}
Le tableau suivant répond au pourquoi, au comment et au où de chaque technologie
retenue. Les \textbf{implémentations concrètes} de ce socle (exécution,
SGBD, bibliothèques front-end, style d'API) et le \textbf{comparatif justifié}
des alternatives évaluées sont détaillés à l'annexe~\ref{chap:technologies}.
\begin{xltabular}{\textwidth}{>{\bfseries}p{2.6cm} X p{5.4cm}}
\toprule
\textbf{Technologie} & \textbf{Pourquoi} & \textbf{Où} \\
\midrule
\endhead
Spring Boot 3 & Socle applicatif standard du marché, conteneur d'injection et autoconfiguration & Organisation générale de l'application \\
Spring MVC & Traiter les requêtes HTTP et orchestrer les traitements & Couche de contrôle \\
React + TypeScript & Produire une interface riche, typée et testable & Couche de présentation \\
Spring Data JPA & Accéder aux données sans écrire le CRUD à la main & Couche de persistance \\
jOOQ & Écrire en SQL typé les requêtes analytiques et PostGIS que JPA exprime mal & Tableaux de bord et géotraitements \\
OpenAPI 3 & Publier un contrat machine et générer les clients tiers & Interface d'interopérabilité \\
Ressources i18n & Externaliser les textes pour la traduction hors-ligne & Module d'internationalisation \\
Profils Spring & Paramétrer l'infrastructure par environnement sans reconstruire l'image & Configuration technique \\
ConfigZumm.ini & Laisser l'exploitant régler ses seuils métier sans intervention technique & Configuration métier \\
Keycloak / OIDC & Authentifier et autoriser sans gérer de mots de passe ni recoder le RBAC & Accès au système \\
X.509 et TLS 1.3 & Chiffrer et authentifier les échanges & Communications inter-composants \\
\bottomrule
\end{xltabular}

\section{Packaging et déploiement}
\begin{itemize}
  \item \textbf{Packaging} : l'application serveur est empaquetée en \textbf{archive Java exécutable} (JAR autoporteur, serveur Tomcat embarqué), puis en \textbf{image conteneur} incluant les fichiers de ressources d'internationalisation. Le client React est construit en fichiers statiques servis par le proxy inverse.
  \item \textbf{Exécution} : déploiement des conteneurs orchestrés, sans serveur d'applications externe à administrer.
  \item \textbf{Base de données} : \textbf{PostgreSQL} étendu par \textbf{PostGIS} (géolocalisation) et \textbf{TimescaleDB} (séries temporelles de mesures).
  \item \textbf{Sécurité} : terminaison \textbf{TLS} au proxy inverse avec certificats X.509, aucun port applicatif ni port d'administration publié directement, configuration du domaine d'identité Keycloak.
  \item \textbf{Configuration} : l'infrastructure est réglée par profil Spring et variables d'environnement ; les seuils, unités et modules métier restent ajustés dans \texttt{ConfigZumm.ini}, monté en volume et rechargeable sans redéploiement de l'image.
\end{itemize}
La répartition physique est illustrée par le diagramme de déploiement (figure \ref{fig:deploiement}).
